\documentclass[UTF8,12pt]{ctexart}
\usepackage[a4paper, top=2.5cm, bottom=2.5cm, left=3cm, right=3cm]{geometry}
\usepackage{graphicx}
\usepackage{float}    
\usepackage{lastpage}
\usepackage{color}
\usepackage{wrapfig}
\usepackage{tikz}%
\usepackage{harpoon}
\usepackage{pgfplots}
\usepackage{tikz-3dplot}
\usetikzlibrary{shapes.geometric, arrows}
\usepackage{ulem}
\usepackage{titlesec}
\usepackage{colortbl}
\usepackage{amsmath}

\usepackage{physics}
\usepackage{listings}
\usepackage{makecell}
\usepackage{indentfirst}
\usepackage{fancyhdr} 
\usepackage{setspace}	
\usepackage{bm}    
\usepackage{siunitx}	
\usepackage{amsfonts,amssymb,times}
\usepackage{txfonts}
\usepackage{enumerate}	
\usepackage{tkz-euclide}
\usetikzlibrary{automata,positioning}
\usepackage{background}
\usepackage[firstpage]{draftwatermark}
\newcommand{\diff}{\mathop{}\!\mathrm{d}}  % 用于积分和微分中的d
\SetWatermarkText{}	
\SetWatermarkLightness{0.85}
\SetWatermarkScale{0.4}
\SetBgContents{}	
\lstset{language=C,keywordstyle=\color{red},showstringspaces=false,rulesepcolor=\color{green}}
\graphicspath{{images/}{figures/}}

% 在每个页面右下角添加"powered by LaTeX"字样
\usepackage{everypage}
\AddEverypageHook{
	\begin{tikzpicture}[remember picture,overlay]
		\node[anchor=south east, 
		xshift=0pt, 
		yshift=0pt, 
		text=gray, 
		opacity=0.5,
		font=\footnotesize\itshape] 
		at (current page.south east) 
		{Editor:HKTK;powered by \LaTeX};
	\end{tikzpicture}
}

% 设置行距为1.5倍
\renewcommand{\baselinestretch}{1.5}
\setlength{\parskip}{0.3em}  % 减小段落间距

% 增大数学公式字号
\usepackage{mathptmx}  % 使用Times字体的数学符号
\everymath{\displaystyle}  % 所有行内公式使用展示样式

% 或者使用以下方式手动调整数学字体大小
\DeclareMathSizes{12}{14}{10}{8}  % 正文12pt时，数学字体14pt，上下标10pt，次下标8pt

\newsavebox{\zdx}	
\newcommand{\putzdx}{
	\marginpar{
		\vspace{-1.5cm}
		\rotatebox[origin=c]{90}{
			\usebox{\zdx}
		}
	}
}

\newcommand{\me}{\mathrm{e}}  
\newcommand{\mi}{\mathrm{i}}
\newcommand{\mj}{\mathrm{j}}
\newcommand{\dif}{\mathrm{d}}
\newcommand{\bs}{\boldsymbol}
\newcommand{\ds}{\displaystyle}

%选择题格式
\newcommand{\fourch}[4]{\vspace{2pt}\\\begin{tabular}{*{4}{@{}p{3.5cm}}}A.~#1 & B.~#2 & C.~#3 & D.~#4\end{tabular}} 
\newcommand{\twoch}[4]{\vspace{2pt}\\\begin{tabular}{*{2}{@{}p{7cm}}}A.~#1 & B.~#2\end{tabular}\\\begin{tabular}{*{2}{@{}p{7cm}}}C.~#3 & D.~#4\end{tabular}}  
\newcommand{\onech}[4]{\\A.~#1\vspace{2pt}\\ B.~#2\vspace{2pt}\\ C.~#3\vspace{2pt}\\ D.~#4}  

\renewcommand{\headrulewidth}{0pt}
\pagestyle{fancy}

% 设置章节标题格式 - 减少间距
\titleformat{\section}{\normalfont\Large\bfseries}{\thesection}{0.5em}{}
\titleformat{\subsection}{\normalfont\large\bfseries}{\thesubsection}{0.5em}{}
\titlespacing*{\section}{0pt}{1.5ex plus 0.5ex minus .2ex}{0.8ex plus .2ex}
\titlespacing*{\subsection}{0pt}{1.2ex plus 0.5ex minus .2ex}{0.6ex plus .2ex}

% 自定义大题标题格式（不换行）
\newcommand{\bigtitle}[1]{\noindent\large\bfseries #1\vspace{2pt}}

% 调整enumerate环境间距
\usepackage{enumitem}
\setlist[enumerate]{topsep=0pt, partopsep=0pt, parsep=2pt, itemsep=2pt, leftmargin=2em}
\setlist[enumerate,1]{label=\arabic*.}
\setlist[enumerate,2]{label=(\arabic*)}

\begin{document} 
	
	\fancyhf{} 
	
	\renewcommand{\hrulefill}{%
		\leavevmode\leaders\hrule height 1.25pt\hfill\kern0pt }
	\newcommand{\blanks}[1]{\underline{\hspace{1cm}}}
	\newcommand{\twowh}[2]{\begin{tabular}{*{2}{@{}p{8cm}}}(1)~#1 & (2)~#2\end{tabular}}
	
	\makeatletter
	\def\lastfox@putlabel{%
		\immediate\write\@auxout{%
			\string\newlabel{LastFox}{{\@arabic\c@fox}{}{}{}{}}}}
	\AtEndDocument{\lastfox@putlabel}
	\makeatother
	\newcounter{foo}
	\newcounter{fox}
	\addtocounter{foo}{1}
	\addtocounter{fox}{0} 
	\fancyfoot[CE,CO]
	
	\vspace{10pt}
	\vspace{16pt}
	
	\begin{spacing}{0}
		\begin{center}
			\textbf{\huge \kaishu 中南大学考试试卷}\\[10pt]
			\textbf{\large  \kaishu 2024\~{}2025学年\underline{一}学期\underline{高等数学A（一）}}\\[10pt]
			\kaishu(时间：2025年1月13日，星期一，10:00—11:40，共计：100分钟)\\[10pt]
			\underline{80}学时，\underline{5}学分，\underline{闭卷}，总分100分，占总评成绩\underline{60}\%
		\end{center}
	\end{spacing}
	
	\subsubsection*{一、填空题（每小题4分，共16分）}
	\begin{enumerate}
		\item 设参数方程 $\begin{cases} x = t - \ln(1 + t) \\ y = t^3 + t^2 \end{cases}$ 确定了函数 $y = y(x)$，则 $\frac{\diff ^2 y}{\diff x^2} |_{t=1} = \underline{\hspace{1cm}}$.
		
		\item 设 $f(x) = \frac{1}{1+x^2} + x^3 \int_0^1 f(x) \diff x$，则 $\int_0^1 f(x) \diff x = \underline{\hspace{2cm}}$.
		
		\item 设 $f(x) = \begin{cases} -1, & -\pi < x \leq 0 \\ 1 + x^2, & 0 < x \leq \pi \end{cases}$，则 $f(x)$ 的以 $2\pi$ 为周期的 Fourier 级数在 $x=\pi$ 处收敛于 \underline{\hspace{2cm}}.
		
		\item $\lim\limits_{n \to \infty} \frac{1}{n}\left[ \sin \frac{\pi}{n} + \sin \frac{2\pi}{n} + \cdots + \sin \frac{(n-1)\pi}{n} \right] = \underline{\hspace{2cm}}$.
	\end{enumerate}
	
	\subsubsection*{二、选择题（每小题4分，共16分）}
	\begin{enumerate}
		\item 设函数 $f(x) = \frac{1 - 2e^\frac{1}{x}}{1 + e^\frac{1}{x}} \arctan x$，则 $x=0$ 是 $f(x)$ 的（ ）
		\fourch{跳跃间断点}{可去间断点}{无穷间断点}{振荡间断点}
		
		\item 设 $f(x)$ 可导，$F(x) = f(x)(1 + |\sin x|)$，则 $f(0)=0$ 是 $F(x)$ 在 $x=0$ 处可导的（ ）
		\twoch{充分必要条件}{充分条件但非必要条件}{必要条件但非充分条件}{既非充分条件又非必要条件}
		
		\item 下列广义积分发散的是（ ）
		\fourch{$\int_{-1}^1 \frac{\diff x}{\sin x}$}{$\int_{-1}^1 \frac{\diff x}{\sqrt{1-x^2}}$}{$\int_0^{+\infty} x e^{-x^2} \diff x$}{$\int_2^{+\infty} \frac{\diff x}{x \ln^2 x}$}
		
		\item 设 $u_n > 0 (n = 1, 2, \ldots)$ 且 $\lim\limits_{n \to \infty} \frac{n}{u_n} = 1$，则级数 $\sum\limits_{n=1}^{\infty} (-1)^{n+1} \left( \frac{1}{u_n} + \frac{1}{u_{n+1}} \right)$（ ）
		\fourch{绝对收敛}{条件收敛}{发散}{收敛性不确定}
	\end{enumerate}
	
	\subsubsection*{三、（本题10分）}
	确定常数 $A, B, C$ 的值，使得当 $x \to 0$ 时，$f(x) = 1 + Ax - e^x (1 + Bx + Cx^2)$ 是比 $x^3$ 高阶的无穷小.
	
	\subsubsection*{四、（本题12分）}
	在第一象限内，求曲线 $y = 1 - x^2$ 上的一点，使该点处的切线与所给的曲线及两坐标轴所围成的图形面积最小，并求此最小面积.
	
	\subsubsection*{五、（本题10分）}
	讨论函数 $f(x) = (2x - 1) \sqrt[3]{(1 - x)^2}$ 的单调性，并求其极值.
	
	\subsubsection*{六、（每小题8分，共16分）}
	\twowh{求极限 $\lim\limits_{x \to 0} \left[ \frac{(1 + x)^{\frac{1}{x}}}{e} \right]^\frac{1}{x}$}{求积分 $\int_0^1 \frac{\ln(1 + x)}{(2 - x)^2}\diff x$}
	
	
	\subsubsection*{七、（本题12分）}
	设函数 $f(x) = \arctan \frac{1-2x}{1+2x}$ 试将 $f(x)$ 展开成 $x$ 的幂级数，并求级数 $\sum\limits_{n=1}^{\infty} \frac{(-1)^n}{2n+1}$ 的和.
	
	\subsubsection*{八、（本题8分）}
	设函数 $f(x)$ 在 $(-\infty, +\infty)$ 是以 $2T$ 为周期的连续函数，证明对任意实数 $x_0$，方程 $f(x) = f(x+T)$ 在区间 $[x_0 - \frac{T}{2}, x_0 + \frac{T}{2}]$ 上至少有一个根.
	
	\newpage
	\begin{spacing}{0}
		\begin{center}
			\textbf{\huge \kaishu 中南大学考试试卷}\\[10pt]
			\textbf{\large  \kaishu 2023\~{}2024学年\underline{一}学期\underline{高等数学A（一）}}\\[10pt]
			\kaishu(时间：2024年1月15日，星期一，10:00—11:40，共计：100分钟)\\[10pt]
			\underline{80}学时，\underline{5}学分，\underline{闭卷}，总分100分，占总评成绩\underline{60}\%
		\end{center}
	\end{spacing}
	\subsubsection*{一、选择题（每小题4分，共16分）}
	\begin{enumerate}
		\item 若 $\lim\limits_{x \to 1} \frac{f(x)}{\ln x} = 1$，则 $\lim\limits_{x \to 1} f(x) $（ ）.
		\twoch{等于1}{等于0}{不存在}{不能确定是否存在}
		
		\item 广义积分 $\int_{0}^{\infty} \frac{1}{x^2 + 2x + 2} \diff x $（ ）.
		\fourch{发散}{等于$\pi$}{等于$\frac{\pi}{2}$}{等于$\frac{\pi}{4}$}
		
		\item 下列四个数项级数，其中收敛的是（ ）.
		\fourch{$\sum\limits_{n=1}^{\infty} \frac{1}{n^{1+\frac{1}{n}}}$}{$\sum\limits_{n=1}^{\infty} \frac{(-1)^n n^{n+1}}{(n+1)!}$}{$\sum\limits_{n=1}^{\infty} \frac{(-1)^n}{\sqrt{n+1}}$}{$\sum\limits_{n=1}^{\infty} \frac{1}{\sqrt{n^2+n}}$}
		
		\item 设函数 $f(x) = \begin{cases} -1, & -\pi < x \leq 0 \\ x^2 + 1, & 0 < x \leq \pi \end{cases}$，周期为$2\pi$，则其傅里叶级数在$x = \pi$处收敛于（ ）.
		\fourch{$-1$}{$f(x)$}{$\pi^2 + 1$}{$\frac{\pi^2}{2}$}
	\end{enumerate}
	
	\subsubsection*{二、填空题（每小题4分，共16分）}
	\begin{enumerate}
		\item 曲线 $y = \frac{2x^2 - x}{1 + x^2}$ 的水平渐近线方程为\underline{\hspace{2cm}}.
		
		\item 若函数 $f(x) = \begin{cases} \frac{\cos(2x) - \cos(3x)}{x^2}, & x \neq 0 \\ A, & x = 0 \end{cases}$ 在 $x = 0$ 处连续，则常数 $A =$ \underline{\hspace{2cm}}.
		
		\item 设函数由方程 $1 + xy = e^y$ 确定，则 $y''|_{x=0} =$ \underline{\hspace{2cm}}.
		
		\item 不定积分 $\int \frac{1}{x(1 + 2\ln x)} \diff x =$ \underline{\hspace{4cm}}.
	\end{enumerate}
	
	\subsubsection*{三、（本题 16 分，每小题 8 分）}
	\begin{enumerate}
		\item 计算定积分 $\int_{-1}^{1} \frac{2x^2 + x \cos x}{1 + \sqrt{1 - x^2}} \diff x$.
		
		\item 已知参数方程 $\begin{cases} x = \cos t + t \sin t \\ y = \sin t - t \cos t \end{cases}$ 确定了函数 $y = y(x)$，求 $\frac{\diff y}{\diff x}$，$\frac{\diff ^2 y}{\diff x^2}$.
	\end{enumerate}
	
	\subsubsection*{四、（本题 10 分）}
	当 $x \to 0$ 时，函数 $f(x) = ax + bx^2 + \ln(1 + x)$ 与函数 $g(x) = e^{x^2} - \cos x$ 是等价无穷小，求 $a, b$ 的值.
	
	\subsubsection*{五、（本题 10 分）}
	将函数 $f(x) = \frac{1}{x}$ 展开成 $x-3$ 的幂级数，并求出收敛域.
	
	\subsubsection*{六、（本题 10 分）}
	求抛物线 $y = -x^2 + 4x$ 与此抛物线上点 (0,0) 和点 (4,0) 处切线所围成的平面图形的面积.
	
	\subsubsection*{七、（本题 12 分）}
	设函数 $f(x)$ 连续，且满足 $\int_{0}^{x}(x-t)f(t)\diff t = x(x-2)e^{x} + 2x$，
	1. 求函数 $f(x)$ 的表达式；2. 求函数 $f(x)$ 的单调区间与极值.
	
	\subsubsection*{八、（本题 10 分）}
	设函数 $f(x)$ 在 $[a,b]$ 上连续，在 $(a,b)$ 内可导，且 $f'(x) \neq 0$，证明：存在 $\xi \in (a,b), \eta \in (a,b)$，使得 $\frac{f'(\xi)}{f'(\eta)} = \frac{e^b - e^a}{b-a} e^{-\eta}$ 成立.
	
	\newpage
	\begin{spacing}{0}
		\begin{center}
			\textbf{\huge \kaishu 中南大学考试试卷}\\[10pt]
			\textbf{\large  \kaishu 2022\~{}2023学年\underline{一}学期\underline{高等数学A（一）}}\\[10pt]
			\kaishu(时间：2023年2月27日，星期一，10:00—11:40，共计：100分钟)\\[10pt]
			\underline{80}学时，\underline{5}学分，\underline{闭卷}，总分100分，占总评成绩\underline{60}\%
		\end{center}
	\end{spacing}
	
	\subsubsection*{一、填空题（每小题4分，共24分）}
	\begin{enumerate}
		\item $\lim\limits_{x \to 0} \frac{(1+2x)^{3x}-1}{x^2} = \underline{\hspace{2cm}}$.
		
		\item 设 $f(x) = x^2 \ln(1-x)$，则 $f^{(5)}(0) = \underline{\hspace{2cm}}$.
		
		\item 若 $x \ln x$ 为函数 $f(x)$ 的一个原函数，则 $\int xf'(x)\diff x = \underline{\hspace{2cm}}$.
		
		\item 设 $f(x)$ 为连续函数，且 $\int_0^x f(t)\diff t = x$，则 $f(8) = \underline{\hspace{2cm}}$.
		
		\item 曲线 $\begin{cases} x = \arctan t \\ y = \frac{1}{2} \ln(1+t^2) \end{cases}$ 上自 $t=0$ 到 $t=1$ 的一段弧长为 \underline{\hspace{2cm}}.
		
		\item 设 $f(x)$ 是周期为 $2\pi$ 的函数，它在一个周期上的表达式为 
		\begin{center}$f(x) = \begin{cases} \frac{1}{\pi} (x+\pi)^2, & -\pi \leq x < 0 \\ \frac{1}{\pi} x^2, & 0 \leq x < \pi \end{cases}$，
		\end{center}\normalsize记 $f(x)$ 展开的傅里叶级数的和函数为 $S(x)$，则 $S(2\pi) = \underline{\hspace{2cm}}$.
	\end{enumerate}
	
	\subsubsection*{二、（本题10分）}讨论函数 $f(x) = \frac{x+1}{\ln |x|}$ 的连续性，写出其连续区间，并指出其间断点的类型.
	
	\subsubsection*{三、（本题 10 分）求下列函数的导数与微分}
	\begin{enumerate}
		\item （5 分）设 $y = \frac{(x)^{\ln x}}{(\ln x)^x}$ ($x > e$)，求 $\diff y$.
		
		\item （5 分）设 $g(x)$ 与 $h(x)$ 是可导函数，求函数 $y = \sqrt{g^2(x) + h^2(x)}$ 的导数（其中 $g(x)$ 与 $h(x)$ 不同时为零）.
	\end{enumerate}
	
	\subsubsection*{四、（本题 12 分）}设由不等式 $0 \leq y \leq \frac{1}{2x+1}$，$0 \leq x \leq 1$ 确定一个平面区域，
	\begin{enumerate}
		\item 求此平面区域的面积；
		\item 求此平面区域绕 y 轴旋转一周得到的立体的体积.
	\end{enumerate}
	
	\subsubsection*{五、（本题 12 分）计算下列积分}
	\begin{enumerate}
		\item （6 分）$\int \frac{3x-2}{x^2-4x+5} \diff x$.
		
		\item （6 分）$\int_0^{\pi/3} \frac{1}{3\sin^2 x + 4\cos^2 x} \diff x$.
	\end{enumerate}
	
	\subsubsection*{六、（本题 10 分）}
	将函数 $f(x) = \frac{x}{x^2 - 3x + 2}$ 展开成麦克劳林级数，并写出其收敛域.
	
	\subsubsection*{七、（本题 16 分）}
	设函数 $f(x) = \frac{x^2 + 3}{x - 1}$，求
	\begin{enumerate}
		\item 函数 $f(x)$ 的单调区间，凹凸区间及极值；
		\item 函数 $f(x)$ 的所有渐近线.
	\end{enumerate}
	
	\subsubsection*{八、（本题 6 分）}
	设 $f(x), g(x)$ 在 $[a, b]$ 上连续，在 $(a, b)$ 内可导，且 $f(x_1) = f(x_2) = 0, (x_1 < x_2, \text{ 且 } x_1, x_2 \in (a, b))$，证明：在 $(x_1, x_2)$ 内至少存在一点 $\xi$，使 $f'(\xi) + f(\xi)g'(\xi) = 0$.
	
	\newpage
	\begin{spacing}{0}
		\begin{center}
			\textbf{\huge \kaishu 中南大学考试试卷}\\[10pt]
			\textbf{\large  \kaishu 2021\~{}2022学年\underline{一}学期\underline{高等数学A（一）}}\\[10pt]
			\kaishu(时间：2022年1月3日，星期一，10:00—11:40，共计：100分钟)\\[10pt]
			\underline{80}学时，\underline{5}学分，\underline{闭卷}，总分100分，占总评成绩\underline{60}\%
		\end{center}
	\end{spacing}
	
	\subsubsection*{一、填空题（每小题4分，共24分）}
	\begin{enumerate}
		\item 当 $x \to 0$ 时，$\sqrt[3]{1 + ax^2 }- 1$ 与 $\cos x - 1$ 是等价无穷小，则 $a = \underline{\hspace{2cm}}$.
		
		\item 曲线 $\begin{cases} x = t^2 \\ y = 3t + t^3 & \end{cases} ,t \in [0, +\infty)$ 的拐点是 \underline{\hspace{2cm}}.
		
		\item 极限 $\lim\limits_{n \to +\infty} \left( \frac{1}{n+1} + \frac{1}{n+2} + \cdots + \frac{1}{n+n} \right) = \underline{\hspace{2cm}}$.
		
		\item 函数 $f(x) = \frac{x}{x+1} \sqrt{1 + \frac{1}{x^2}}$ 的无穷间断点是 \underline{\hspace{2cm}}.
		
		\item 曲线 $y = \int_0^x \tan t \, \diff t \ (0 \leq x \leq \frac{\pi}{4})$ 的弧长是 \underline{\hspace{2cm}}.
		
		\item 将以$2\pi$为周期的函数 $f(x) = \frac{\pi - x}{2} \ (0 < x < 2\pi)$ 展开成 Fourier 级数，记级数的和函数为 $S(x)$，则 $S(14) = \underline{\hspace{2cm}}$.
	\end{enumerate}
	
	\subsubsection*{二、（本题10分）}
	设函数 $f(x)$ 定义如下： $f(x) = \begin{cases} \frac{1}{\ln{(1+x)}} - \frac{1}{x}, & x > 0, \\ a + bx, & x \leq 0. \end{cases}$ 确定常数 $a, b$，使得 $f(x)$ 在点 $x = 0$ 可导.
	
	\subsubsection*{三、（本题 8 分）}
	将半径为 $R$，夹角为 $\theta (0 < \theta < 2\pi)$ 的扇形围成一个无底的圆锥体，问夹角 $\theta$ 取多大时，圆锥体的体积最大？
	
	\subsubsection*{四、（本题 14 分，每小题 7 分）}
	\begin{enumerate}
		\item 计算不定积分 $\int x f'(2x) \diff x$，其中 $f(x)$ 的原函数为 $\frac{\sin x}{x}$.
		
		\item 计算定积分 $\int_{-\pi/2}^{\pi/2} \frac{\sin^2 x}{1+e^{-x}} \diff x$.
	\end{enumerate}
	
	\subsubsection*{五、（本题 12 分）}
	设曲线 $y = x^2 - 2x$，$y = 0$，$x = 1$，$x = 3$ 所围成的平面图形为 $D$，求
	\begin{enumerate}
		\item 平面图形 $D$ 的面积 $S$；
		\item 该平面图形 $D$ 绕 $y$ 轴旋转一周所得旋转体的体积 $V$.
	\end{enumerate}
	
	\subsubsection*{六、（本题 10 分）}
	已知方程 $\begin{cases} x = 3t^2 + 2t + 3 \\ e^{y} \sin t - y + 1 = 0 \end{cases}$ 确定了 $y$ 是 $x$ 的函数，求 $\frac{\diff y}{\diff x} \bigg|_{t=0}, \frac{\diff ^2 y}{\diff x^2} \bigg|_{t=0}$.
	
	\subsubsection*{七、（本题 16 分，每小题 8 分）}
	\begin{enumerate}
		\item 判定级数 $-\frac{1}{1+a} + \frac{1}{2+a} - \frac{1}{3+a} + \frac{1}{4+a} - \cdots$ (a 为自然数) 是否收敛？若收敛，是绝对收敛还是条件收敛？
		
		\item 求级数 $\sum\limits_{n=1}^{\infty} (-1)^{n-1} \frac{x^{2n}}{2n-1}$ 的收敛域及和函数.
	\end{enumerate}
	
	\subsubsection*{八、（本题 6 分）}
	设函数 $f(x)$ 在 $[0,1]$ 上连续，在 $(0,1)$ 可导，且 $f(0) = f(1) = 0$，$f\left(\frac{1}{2}\right) = 1$, 证明至少存在一点 $\xi \in (0,1)$，使得 $f'(\xi) = 1$.
	
	\newpage
	\begin{spacing}{0}
		\begin{center}
			\textbf{\huge \kaishu 中南大学考试试卷}\\[10pt]
			\textbf{\large  \kaishu 2020\~{}2021学年\underline{一}学期\underline{高等数学A（一）}}\\[10pt]
			\kaishu(时间：2021年1月22日，星期五，10:00—11:40，共计：100分钟)\\[10pt]
			\underline{80}学时，\underline{5}学分，\underline{闭卷}，总分100分，占总评成绩\underline{60}\%
		\end{center}
	\end{spacing}
	\subsubsection*{一、填空题(每小题4分，共24分)}
	\begin{enumerate}
		\item 设函数 $y = y(x)$ 由参数方程 $\begin{cases} x = t + e^t \\ y = \sin t \end{cases}$ 确定，则 $\frac{\diff ^2 y}{\diff x^2} \bigg|_{t=0} = \underline{\hspace{2cm}}$.
		
		\item 函数 $f(x) = x^2 \cdot 2^x$ 在 $x = 0$ 处的 n 阶导数 $f^{(n)}(0) = \underline{\hspace{2cm}}$.
		
		\item 曲线 $y = \frac{x^2}{2x+1}$ 的斜渐近线方程为 \underline{\hspace{4cm}}.
		
		\item $\lim\limits_{n \to \infty} n \left( \frac{1}{1+n^2} + \frac{1}{2^2+n^2} + \cdots + \frac{1}{n^2+n^2} \right) = \underline{\hspace{2cm}}$.
		
		\item 当 $0 \leq \theta \leq \pi$ 时，对数螺线 $r = e^{\theta}$ 的弧长为 \underline{\hspace{2cm}}.
		
		\item 设函数 $f(x) = \pi x + x^2 (-\pi < x < \pi)$ 的 Fourier 展开式为 $\frac{a_0}{2} + \sum\limits_{n=1}^{\infty} (a_n \cos nx + b_n \sin nx)$，则其中系数 $b_3 = \underline{\hspace{2cm}}$.
	\end{enumerate}
	
	\subsubsection*{二、（满分10分）}
	设函数 $f(x) = x + a \ln (1 + x) + b \sin x, g(x) = kx^3$，若 $f(x)$ 与 $g(x)$ 在 $x \to 0$ 时是等价无穷小，求 $a, b, k$ 的值.
	
	\subsubsection*{三、（满分16分）}
	设 $f(x) = \begin{cases} \frac{g(x)-e^{-x}}{x}, & x \neq 0 \\ 0, & x = 0 \end{cases}$ 其中 $g(x)$ 有二阶连续导数，且 $g(0) = 1$，$g'(0) = -1$.
	\begin{enumerate}
		\item 求 $f'(x)$.
		\item 讨论 $f'(x)$ 在 $(-\infty, +\infty)$ 上的连续性.
	\end{enumerate}
	
	\subsubsection*{四、（满分10分）}
	设函数 $y = y(x)$ 由方程 $y \ln y - x + y = 0$ 确定，试判断曲线 $y = y(x)$ 在 (1,1) 附近的凹凸性.
	
	\subsubsection*{五、（满分8分）}
	设 $f(\sin^2 x) = \frac{x}{\sin x}$，求 $\int \frac{\sqrt{x}}{\sqrt{1-x}} f(x) \diff x$.
	
	\subsubsection*{六、（满分16分）}
	设 $D_1$ 是由抛物线 $y = 2x^2$ 和直线 $x = a, x = 2$ 及 $y = 0$ 所围成的平面区域，$D_2$ 是由抛物线 $y = 2x^2$ 和直线 $x = a, y = 0$ 所围成的平面区域，其中 $0 < a < 2$.
	\begin{enumerate}
		\item 求 $D_1$ 绕 $x$ 轴旋转而成的旋转体体积 $V_1$，$D_2$ 绕 $y$ 轴旋转而成的旋转体体积 $V_2$；
		\item 当 $a$ 为何值时，$V = V_1 + V_2$ 取得最大值？求此最大值.
	\end{enumerate}
	
	\subsubsection*{七、（满分10分）}
	将函数 $y = \ln(1 - x - 2x^2)$ 展开成 $x$ 的幂级数，并指出其收敛区间.
	
	\subsubsection*{八、（满分6分）}
	设 $f(x)$ 在 $[a, b]$ 上连续，且在 $(a, b)$ 内可导，且 $f'(x) \leq 0$，记 $F(x) = \frac{1}{x-a} \int_a^x f(t) \diff t$，证明在 $(a, b)$ 内 $F'(x) \leq 0$.
	
	\newpage
	\begin{spacing}{0}
		\begin{center}
			\textbf{\huge \kaishu 中南大学考试试卷}\\[10pt]
			\textbf{\large  \kaishu 2019\~{}2020学年\underline{一}学期\underline{高等数学A（一）}}\\[10pt]
			\kaishu(时间：2020年1月9日，星期四，10:00—11:40，共计：100分钟)\\[10pt]
			\underline{80}学时，\underline{5}学分，\underline{闭卷}，总分100分，占总评成绩\underline{60}\%
		\end{center}
	\end{spacing}
	
	\subsubsection*{一、填空题(每小题4分，共24分)}
	\begin{enumerate}
		\item 已知 $x \to 0$，$(1 - ax^2)^{\frac{1}{4}} - 1$ 与 $x \sin x$ 是等价无穷小，求 $a = \underline{\hspace{2cm}}$.
		
		\item 求极限： $\lim\limits_{n \to \infty} \frac{1 + 2! + \cdots + n!}{n!} = \underline{\hspace{2cm}}$.
		
		\item 求 $f(x) = 2^x$ 的麦克劳林的 $x^n$ 项前的系数 \underline{\hspace{2cm}}.
		
		\item 函数 $y(x)$ 由 $\begin{cases} x = t^3 + 3t + 1 \\ y = t^3 - 3t + 1 \end{cases}$ 确定，若曲线$y(x)$ 向上凸，则$x$ 的取值范围 \underline{\hspace{2cm}}.
		
		\item 设 $a \geq 0$，求 $\int_0^{+a} \frac{1}{(1+x^2)(1+x^a)} \diff x = \underline{\hspace{2cm}}$.
		
		\item 设$S(x)$ 是周期为 $2\pi$ 函数 $f(x)$ 的Fourier级数的和函数，$f(x)$ 在一个周期的表达式 $f(x) = \begin{cases} 0, & 2 < |x| < \pi \\ x, & |x| \leq 2 \end{cases}$，求 $S(6) = \underline{\hspace{2cm}}$.
	\end{enumerate}
	
	\subsubsection*{二、（本题10分）}
	设 $f(x) = \lim\limits_{n \to \infty} \frac{x^{2n-1} + ax^2 + bx}{x^{2n+1}}$ 是 $(-\infty, +\infty)$ 上的连续函数，试求 $a, b$ 的值.
	
	\subsubsection*{三、（本题10分）}
	设 $a > 1, f(t) = a^t - at$ 在 $(-\infty, +\infty)$ 内的驻点为 $t(a)$.问 $a$ 为何值，$t(a)$ 最小？并求出最小值.
	
	\subsubsection*{四、（本题16分，每小题7分）}
	\begin{enumerate}
		\item 设 $f(x) = \int_{x^2}^{x} \frac{\sin(xt)}{t} \, \diff t$，求 $\lim\limits_{x \to 0} \frac{f(x)}{x^2}$.
		
		\item 求不定积分 $\int \frac{e^{\arctan x}}{(1+x^2)^\frac{3}{2}} \, \diff x$.
	\end{enumerate}
	
	\subsubsection*{五、（本题16分，每小题8分）}
	过坐标原点作曲线$y=\ln x$的切线，该切线与曲线$y=\ln x$及$x$轴围成平面图形$D$。
	\\(1)求$D$的面积;
	\\(2)求$D$绕直线$x=e$旋转一周所得旋转体的体积$V$.
	
	\subsubsection*{六、（本题10分）}
	求出级数 $\sum\limits_{n=1}^{\infty} (-1)^{n-1} \left[ 1 + \frac{1}{n(2n-1)} \right] x^{2n}$ 收敛区间与和函数.
	
	\subsubsection*{七、（本题8分）}
	求出 $y = \frac{(x^2 + 2x - 3)e^{\frac{1}{x}}}{(x^2 - 1)\arctan x}$ 的所有渐近线.
	
	\subsubsection*{八、（本题6分）}
	设函数$f(x)$ 在 $[a, b]$ 上连续.在 $(a, b)$ 上可导.且 $f'(x) > 0$.证明存在唯一点 $\xi \in (a, b)$，使 $\int_a^\xi [f(\xi) - f(x)] \, \diff x = 3 \int_\xi^b [f(x) - f(\xi)] \, \diff x, \xi \in (a, b)$.
	
	\newpage
	\begin{spacing}{0}
		\begin{center}
			\textbf{\huge \kaishu 中南大学考试试卷}\\[10pt]
			\textbf{\large  \kaishu 2018\~{}2019学年\underline{一}学期\underline{高等数学A（一）}}\\[10pt]
			\kaishu(时间：2019年1月17日，星期四，10:00—11:40，共计：100分钟)\\[10pt]
			\underline{80}学时，\underline{5}学分，\underline{闭卷}，总分100分，占总评成绩\underline{60}\%
		\end{center}
	\end{spacing}
	
	\subsubsection*{一、填空题(每小题 4 分，总计 24 分)}
	\begin{enumerate}
		\item 若函数 $f(x) = \begin{cases} \frac{1 - \cos \sqrt{x}}{ax}, & x > 0 \\ b, & x \leq 0 \end{cases}$ 在 $x = 0$ 处连续，则 $ab = \underline{\hspace{2cm}}$.
		
		\item 已知 $f'(x) = e^{3x}$，则 $f(x)$ 的反函数的二阶导数 $\frac{\diff^2 x}{\diff y^2} = \underline{\hspace{2cm}}$.
		
		\item 已知函数 $f(x) = \frac{1}{1 + x^2}$，则 $f''(0) = \underline{\hspace{2cm}}$.
		
		\item 已知 $f(x)$ 的一个原函数为 $x^2 \sin x$，则 $\int f'(x) \diff x = \underline{\hspace{2cm}}$.
		
		\item 广义积分 $\int_{e}^{\infty} \frac{1}{x(\ln x)^2} \diff x = \underline{\hspace{2cm}}$.
		
		\item 设函数 $f(x)$ 是以 $2\pi$ 为周期的函数，且 $f(x) = \begin{cases} -1, & -\pi < x \leq 0 \\ 1 + x^2, & 0 < x \leq \pi \end{cases}$，则其傅里叶级数在点 $x = 0$ 处的和为 \underline{\hspace{2cm}}.
	\end{enumerate}
	
	\subsubsection*{二、（共 14 分，其中第 1 小题 7 分，第 2 小题 7 分）}
	\begin{enumerate}
		\item 求极限 $\lim\limits_{x \to 0} \frac{2\sin x + x^2 \cos \frac{1}{x}}{\ln (1 + x)}$.
		
		\item 求函数 $f(x) = \frac{x^2 - 1}{|x| (x + 1)}$ 的间断点，并判断其类型.
	\end{enumerate}
	
	\subsubsection*{三、（10 分）}
	证明 Cauchy 中值定理：如果函数 $f(x), g(x)$ 在 $[a, b]$ 上连续，在 $(a, b)$ 内可导，且 $g'(x) \neq 0, x \in (a, b)$，证明至少存在一点 $\xi \in (a, b)$，使得 $\frac{f(b) - f(a)}{g(b) - g(a)} = \frac{f'(\xi)}{g'(\xi)}$.
	
	\subsubsection*{四、（共 14 分，其中第 1 小题 7 分，第 2 小题 7 分）}
	\begin{enumerate}
		\item 试确定 $a, b, c$ 的值，使三次曲线 $y = ax^3 + bx^2 + cx$ 有拐点 $(1, 2)$，并且在该点处切线的斜率为 1.
		
		\item 设函数 $y = e^{-x} \sqrt{\frac{(x+1)(x^2+2)}{3-x^3}}$，求 $\diff y$.
	\end{enumerate}
	
	\subsubsection*{五、（10 分）}
	设 $D$ 是由曲线 $y = x^{\frac{1}{y}}$，直线 $x = a(a > 0)$ 及 $x$ 轴所围成的平面图形，$V_x$、$V_y$ 分别是 $D$ 绕 $x$ 轴，$y$ 轴所得旋转体的体积，若 $V_y = 10V_x$，求 $a$.
	
	\subsubsection*{六、（6 分）}
	设 $f(x) = \int_{0}^{x^2} \ln(1+t)\diff t$，证明导函数 $f'(x)$ 只有一个零点.
	
	\subsubsection*{七、（共 12 分，其中第 1 小题 6 分，第 2 小题 6 分）}
	\begin{enumerate}
		\item 求不定积分 $\int \frac{1}{\sqrt{4-x^2} \arcsin \frac{x}{2}} \diff x$.
		
		\item $y = [x]$ 表示不超过 $x$ 的最大整数，计算 $\int_{-2}^{10} (x-[x]) \diff x$.
	\end{enumerate}
	
	\subsubsection*{八、（10 分）}
	将函数 $f(x) = \frac{1}{x}$ 展开成 $(x-3)$ 的幂级数，并指出其收敛域.
	
	\newpage
	\begin{spacing}{0}
		\begin{center}
			\textbf{\huge \kaishu 中南大学考试试卷}\\[10pt]
			\textbf{\large  \kaishu 2017\~{}2018学年\underline{一}学期\underline{高等数学A（一）}}\\[10pt]
			\kaishu(时间：2018年1月18日，星期四，10:00—11:40，共计：100分钟)\\[10pt]
			\underline{80}学时，\underline{5}学分，\underline{闭卷}，总分100分，占总评成绩\underline{60}\%
		\end{center}
	\end{spacing}
	
	\subsubsection*{一、填空题(本题24分，每小题3分)}
	\begin{enumerate}
		\item 设 $\lim\limits_{x \to 0} \frac{\int_{0}^{x} \frac{t^2}{\sqrt{a + t^2}} \diff t}{x - \sin x} = 1$，则 $a = \underline{\hspace{2cm}}$.
		
		\item 设 $y = \sin f(x^2)$，其中函数 $f(u)$ 是可微的，则 $\diff y = \underline{\hspace{2cm}}$.
		
		\item 设函数 $f(x) = \begin{cases} x \arctan \frac{1}{x}, & x \neq 0 \\ 2, & x = 0 \end{cases}$，则 $f(x)$ 的间断点为\underline{\hspace{2cm}}，间断点的类型为\underline{\hspace{2cm}}.
		
		\item 设 $f'(\sin^2 x) = \cos^2 x$，则 $\int f(x) \diff x = \underline{\hspace{2cm}}$.
		
		\item $\int_{-\frac{\pi}{2}}^{\frac{\pi}{2}} \frac{\sin^3 x}{(1 + x^2)^{20}} \cos x \diff x = \underline{\hspace{2cm}}$.
		
		\item 星形线 $x = a \cos^3 t, y = a \sin^3 t$ 的全长为\underline{\hspace{2cm}}.
		
		\item 已知 $\int_{-\infty}^{\infty} e^{-x^2} \diff x = \sqrt{\pi}$，则广义积分 $\int_{-\infty}^{\infty} e^{-x^2 - x} \diff x = \underline{\hspace{2cm}}$.
		
		\item 设 $f(x) = 1 - \frac{x}{\pi} (0 \leq x \leq \pi)$ 的余弦级数的和函数为 $S(x)$，则 $S(-3) = \underline{\hspace{2cm}}$.
	\end{enumerate}
	
	\subsubsection*{二、（共16分，其中第1小题8分，第2小题8分）}
	\begin{enumerate}
		\item 求极限 $\lim\limits_{n \to \infty} \left( \frac{n}{n^2 + 1} + \frac{n}{n^2 + 2} + \cdots + \frac{n}{n^2 + n} \right)$.
		
		\item 设函数 $f(x) = \begin{cases} \frac{x^4 \cos \frac{1}{x^2}}{e^x - 1}, & x < 0 \\ a, & x = 0 \\ \frac{\ln(1 + 3x)}{x} + b, & x > 0 \end{cases}$ 在 $(-\infty,+\infty)$ 内连续，求 $a, b$.
	\end{enumerate}
	
	\subsubsection*{三、（10 分）}
	设上证 A 股中某只股票历年来的走势指数曲线为 $C: f(x) = \frac{x^2 - 2x + 2}{x - 1}$，求（1）$f(x)$的增减区间和极值；（2）$f(x)$的四凸区间及拐点；（3）$f(x)$的渐近线.
	
	\subsubsection*{四、（8 分）}
	求幂级数 $\sum\limits_{n=1}^{\infty} \frac{2n+1}{n!} x^{2n}$ 的收敛区间，并在其收敛区间内求其和函数.
	
	\subsubsection*{五、（10 分）}
	求曲线 $y = x^2 - 2x$，$y = 0$，$x = 1$，$x = 3$ 所围成的平面图形的面积 $S$，并求该平面图形绕 $x$ 轴旋转一周所得旋转体的体积 $V$.
	
	\subsubsection*{六、（共 12 分，其中第 1 小题 6 分，第 2 小题 6 分）}
	计算1、$\int \frac{2x^5 + 11x^4 + 28x^3 + 39x^2 + 32x + 13}{2x^3 + 7x^2 + 8x + 3} \diff x$；\\ 2、$\int_{\frac{1}{2}}^{1} \frac{\arcsin \sqrt{x}}{\sqrt{x(1-x)}} \diff x$.
	
	\subsubsection*{七、(8 分)}
	设 $\{a_n\}$ 是单调增加且有界的数列，(其中 $a_n > 0, n = 1, 2, 3, \cdots$)，证明级数 $\sum\limits_{n=1}^{\infty} \left( 1 - \frac{a_n}{a_{n+1}} \right)$ 收敛.
	
	\subsubsection*{八、(共 12 分，其中第 1 小题 6 分，第 2 小题 6 分)}
	\begin{enumerate}
		\item 设 $f(x)$ 在 $[0,1]$ 上连续，在 $(0,1)$ 内可导，且 $f(1) = 0$，证明至少存在一点 $\eta \in (0,1)$，使得 $\eta f'(\eta) + 2f(\eta) = 0$.
		
		\item 设 $\begin{cases} x = \int_0^t \frac{u}{(2-u)} du \\ y = \sum\limits_{n=1}^{\infty} \frac{n}{2^n} t^{n-1} \end{cases}$，求 $\frac{\diff y}{\diff x}, \frac{\diff^2 y}{\diff x^2}$.
	\end{enumerate}
	
	\newpage
	\begin{spacing}{0}
		\begin{center}
			\textbf{\huge \kaishu 中南大学考试试卷}\\[10pt]
			\textbf{\large  \kaishu 2016\~{}2017学年\underline{一}学期\underline{高等数学A（一）}}\\[10pt]
			\kaishu(时间：2017年1月13日，星期五，10:00—11:40，共计：100分钟)\\[10pt]
			\underline{80}学时，\underline{5}学分，\underline{闭卷}，总分100分，占总评成绩\underline{60}\%
		\end{center}
	\end{spacing}
	
	\subsubsection*{一、计算下列各题(本题 35 分，每小题 7 分)}
	\begin{enumerate}
		\item 求极限 $\lim\limits_{x \to \infty} \frac{1}{x} \int_0^{2x} (1 + t^2) e^{t^2 - x^2} \diff t$.
		
		\item 设 $y = y(x)$ 由方程 $xe^{f(y)} = e^y$ 所确定，其中 $f$ 有一阶导数且 $f' \neq 1$，求 $\frac{\diff y}{\diff x}$.
		
		\item 求不定积分 $\int \frac{x+1}{x(1+xe^x)} \diff x$.
		
		\item 判断级数 $\sum\limits_{n=1}^{\infty} \left( \frac{na}{n+1} \right)^n (a > 0)$ 的敛散性.
		
		\item 求定积分 $\int_0^{\pi} \frac{1}{2+\cos x} \diff x$.
	\end{enumerate}
	
	\subsubsection*{二、（本题 10 分）}
	设曲线的极坐标方程为 $r = a \sin 3\theta$，求它在点 $(a, \frac{\pi}{6})$ 处的法线的直角坐标方程.
	
	\subsubsection*{三、（本题 15 分）}
	求幂级数 $\sum\limits_{n=1}^{\infty} \frac{1}{n!(n+2)} x^{n+2}$ 的收敛域，在收敛域内求该级数的和函数 $S(x)$，并求常数项级数 $\sum\limits_{n=1}^{\infty} \frac{n+2}{n!+(n+1)!+(n+2)!}$ 的和.
	
	\subsubsection*{四、（本题 10 分）}
	设由曲线 $y = x^2$，直线 $y = kx(0 < k < 1)$ 及 $x = 1$ 围成的平面图形面积为 $S(k)$，求 $k$ 的值，使 $S(k)$ 为最小.
	
	\subsubsection*{五、（本题 10 分）}
	当 $x \to 0$ 时，$1 - \cos x \cos 2x \cos 3x$ 与 $ax^n$ 是等价无穷小，求常数 $a$ 与 $n$.
	
	\subsubsection*{六、（本题 10 分）}
	设函数 $f(x)$ 在 $[0, \pi]$ 上有二阶连续导数，且 $\int_0^\pi [f(x) + f''(x)] \sin x \diff x = 5, \quad f(\pi) = 2$，求 $f(0)$.
	
	\subsubsection*{七、（本题 10 分）}
	若 $ab > 0$，证明 $ae^b - be^a = (\xi - 1)e^{\xi} (b - a)$，其中 $\xi$ 在 $a, b$ 之间.
	
	\newpage
	\begin{spacing}{0}
		\begin{center}
			\textbf{\huge \kaishu 中南大学考试试卷}\\[10pt]
			\textbf{\large  \kaishu 2015\~{}2016学年\underline{一}学期\underline{高等数学A（一）}}\\[10pt]
			\kaishu(时间：2016年1月14日，星期四，10:00—11:40，共计：100分钟)\\[10pt]
			\underline{80}学时，\underline{5}学分，\underline{闭卷}，总分100分，占总评成绩\underline{60}\%
		\end{center}
	\end{spacing}
	
	\subsubsection*{一、(共 15 分：其中第一小题 10 分，第二小题 5 分)}
	\begin{enumerate}
		\item 讨论函数 $f(x) = \begin{cases} \frac{1}{1-e^{\frac{x}{x-1}}}, & x \neq 1 \\ 1 , & x = 1 \end{cases}$ 的连续性，若存在间断点指出其类型.
		
		\item 已知 $\lim\limits_{x \to 0} \frac{\sqrt{1 + f(x)\sin x} - 1}{e^{2x} - 1} = 3$，求极限 $\lim\limits_{x \to 0} f(x)$.
	\end{enumerate}
	
	\subsubsection*{二、(10 分)}
	设函数 $f(x) = x + a \ln (1 + x) + bx \sin x$，$g(x) = kx^3$，若 $f(x)$ 与 $g(x)$ 在 $x \to 0$ 时是等价无穷小，求 $a, b, k$ 的值.
	
	\subsubsection*{三、（共 15 分：其中第一小题 5 分，第二小题 10 分）}
	\begin{enumerate}
		\item 计算 $\int_{0}^{3} \frac{1}{(x-1)^{\frac{2}{3}}} \, \diff x$.
		
		\item 求由方程 $\int_{y^2}^{0} e^{t} \diff t + \int_{0}^{2x} t^{2} \diff t = 0$ 所确定的隐函数 $y$ 对 $x$ 的导数 $\frac{\diff y}{\diff x}, \frac{\diff^{2} y}{\diff x^{2}}$.
	\end{enumerate}
	
	\subsubsection*{四、（10 分）}
	设 $f(x)$ 在闭区间 $[0,1]$ 上连续，在开区间 $(0,1)$ 内可导，且 $f(1) = 0$，证明至少存在一点 $\xi \in (0,1)$，使 $\xi f'(\xi) + nf(\xi) = 0$，其中 $n$ 为正整数.
	
	\subsubsection*{五、（共15分：其中第一小题7分，第二小题8分）}
	\begin{enumerate}
		\item 计算不定积分 $\int \frac{e^{\arctan x}}{(1 + x^2)^{\frac{3}{2}}} \, \diff x$.
		
		\item 设 $\int_{-\pi}^{\pi} (x - a_0 \cos x - a_0 \sin x)^2 \, \diff x = \min_{\alpha \in \mathbb{R}} \left[ \int_{-\pi}^{\pi} (x - a \cos x - a \sin x)^2 \, \diff x \right]$，求 $a_0$.
	\end{enumerate}
	
	\subsubsection*{六、（10 分）}
	由原点引抛物线 $y = x^2 + 2x + 4$ 的两条切线，设切点分别为 $A, B$，求两切线 $OA, OB$ 与此抛物线所围成的平面图形的面积.
	
	\subsubsection*{七、（共 15 分：其中第一小题 7 分，第二小题 8 分）}
	\begin{enumerate}
		\item 判断级数 $\sum\limits_{n=1}^{\infty} u_n = 1 - \left( \frac{1}{5} \right)^2 + \frac{1}{3} - \left( \frac{1}{5} \right)^4 + \frac{1}{5} - \left( \frac{1}{5} \right)^6 + \cdots$ 的敛散性.
		
		\item 求曲线 $y = f(x) = \frac{(x^2 + 2x - 3)e^{\frac{1}{x}}}{(x^2 - 1)\arctan x}$ 的所有渐近线方程.
	\end{enumerate}
	
	\subsubsection*{八、（10 分）}
	将函数 $f(x) = \frac{x}{2-x-x^2}$ 展开成 $x$ 的幂级数，并求其收敛域.
	
	\newpage
	\begin{spacing}{0}
		\begin{center}
			\textbf{\huge \kaishu 中南大学考试试卷}\\[10pt]
			\textbf{\large  \kaishu 2014\~{}2015学年\underline{一}学期\underline{高等数学A（一）}}\\[10pt]
			\kaishu(时间：2015年1月22日，星期四，10:00—11:40，共计：100分钟)\\[10pt]
			\underline{80}学时，\underline{5}学分，\underline{闭卷}，总分100分，占总评成绩\underline{60}\%
		\end{center}
	\end{spacing}
	
	\subsubsection*{一、填空题(本题15分，每小题3分)}
	\begin{enumerate}
		\item $\lim\limits_{n \to \infty} \frac{1}{n^2} + \frac{1}{(n+1)^2} + \cdots + \frac{1}{(n+n)^2} = \underline{\hspace{2cm}}$.
		
		\item $x=0$ 为函数 $f(x) = \begin{cases} \arctan \frac{1}{x}, & x \neq 0 \\ 0, & x = 0 \end{cases}$ 的 \underline{\hspace{2cm}} 间断点.
		
		\item 设 $f(x) = \frac{x^2}{1-x^2}$，则 $f^{(2015)} (0) = \underline{\hspace{2cm}}$.
		
		\item $\int x \sec^2 x \diff x = \underline{\hspace{2cm}}$.
		
		\item 周期为2的函数 $f(x)$，它在一个周期上的表达式为 $f(x) = x, -1 \leq x < 1$，设它的傅里叶级数的和函数为 $S(x)$，则 $S(1.5) = \underline{\hspace{2cm}}$.
	\end{enumerate}
	
	\subsubsection*{二、选择题(本题15分，每小题3分)}
	\begin{enumerate}
		\item 设函数 $f(x) = (1+x+\cdots+\frac{x^n}{n!})e^{-x}$，则（ ）
		\onech{当 $n$ 为奇数时，$f(x)$ 在 $x=0$ 取极大值}{当 $n$ 为奇数时，$f(x)$ 在 $x=0$ 取极小值}{当 $n$ 为偶数时，$f(x)$ 在 $x=0$ 取极大值}{当 $n$ 为偶数时，$f(x)$ 在 $x=0$ 取极小值}
		
		\item 曲线 $y = \frac{1+e^{-x^2}}{1-e^{-x^2}}$，则该曲线（ ）
		\twoch{没有渐近线}{仅有水平渐近线}{仅有铅直渐近线}{既有水平渐近线又有铅直渐近线}
		
		\item 下列广义积分收敛的是（ ）
		\fourch{$\int_1^{\infty} \frac{\diff x}{\sqrt{x}}$}{$\int_0^1 \frac{\diff x}{\sqrt{x^3}}$}{$\int_{-1}^1 \frac{\diff x}{x^2}$}{$\int_1^{\infty} x e^{-x} \diff x$}
		
		\item 曲线 $y = \ln x$ 及直线 $x = e$，$x$ 轴所围成的图形绕 $y$ 轴旋转形成的旋转体的体积为 $V_y$，则有 $V_y = $（ ）
		\fourch{$\frac{\pi}{2} e^2$}{$\frac{\pi}{2} (e^2 - 1)$}{$\frac{\pi}{2} (e^2 + 1)$}{$\pi e^2$}
		
		\item 下列级数条件收敛的是（ ）
		\twoch{$\sum\limits_{n=1}^{\infty} \frac{(-1)^n n}{2n+10}$}{$\sum\limits_{n=1}^{\infty} \frac{(-1)^{n-1}}{\sqrt{n^3}}$}{$\sum\limits_{n=1}^{\infty} (-1)^{n-1} \left( \frac{1}{2} \right)^n$}{$\sum\limits_{n=1}^{\infty} (-1)^{n-1} \frac{3}{\sqrt{n}}$}
	\end{enumerate}
	
	\subsubsection*{三、（共14分，每小题7分）求下列各极限}
	\begin{enumerate}
		\item 设函数 $f(x) = \frac{2 + e^x}{1 + e^x} + \frac{\sin x}{|x|}$，求极限 $\lim\limits_{x \to 0} f(x)$.
		
		\item $\lim\limits_{x \to 0} \frac{a^x - a^{\sin x}}{x^3}$ $(a > 0)$.
	\end{enumerate}
	
	\subsubsection*{四、（10分）}
	某房地产公司有50套公寓要出租，当租金定为每月180元时，公寓会全部租出去.当租金每月增加10元时，就有一套公寓租不出去，而租出去的房子每月需花费20元的整修维护费.试问房租定为多少可获得最大收入？
	
	\subsubsection*{五、（共16分，每小题8分）求解导数题}
	\begin{enumerate}
		\item 设曲线方程为 $\begin{cases} x = 2t + 3 + \arctan t \\ y = 2 - 3t + \ln(1 + t^2) \end{cases}$，求该曲线在 $x = 3$ 处的切线方程.
		
		\item 设函数 $y = y(x)$ 由方程 $2y - x = (x - y)\ln(x - y)$ 确定，求 $\frac{\diff ^2 y}{\diff x^2}$.
	\end{enumerate}
	
	\subsubsection*{六、（9分）}
	将函数 $y = \ln(1 - x - 2x^2)$ 展开成 $x$ 的幂级数，并指出其收敛域.
	
	\subsubsection*{七、（共14分，每小题7分）求解下列各题}
	\begin{enumerate}
		\item 已知 $f'(x) = \begin{cases} x^2, & x \geq 0 \\ \sin x, & x < 0 \end{cases}$，求 $f(x)$.
		
		\item 求星形线 $\begin{cases} x = a \cos^3 t \\ y = a \sin^3 t \end{cases}$，$(0 \leq t \leq 2\pi, a > 0)$ 的全长.
	\end{enumerate}
	
	\subsubsection*{八、（7分）}
	设 $f(x), g(x)$ 在 $[a, b]$ 上连续，且 $f(x) > 0, g(x) > 0$，证明：$\lim\limits_{n \to \infty} \int_a^b g(x) \sqrt[n]{f(x)} \diff x = \int_a^b g(x) \diff x$.
	
	\newpage
	\begin{spacing}{0}
		\begin{center}
			\textbf{\huge \kaishu 中南大学考试试卷}\\[10pt]
			\textbf{\large  \kaishu 2013\~{}2014学年\underline{一}学期\underline{高等数学A（一）}}\\[10pt]
			\kaishu(时间：2014年1月16日，星期四，10:00—11:40，共计：100分钟)\\[10pt]
			\underline{80}学时，\underline{5}学分，\underline{闭卷}，总分100分，占总评成绩\underline{60}\%
		\end{center}
	\end{spacing}
	
	\subsubsection*{一、填空题(本题15分，每小题3分)}
	\begin{enumerate}
		\item 当 $x \to 0$ 时，$\sqrt[5]{1 + ax} - 1$ 与 $\cos \sqrt{x} - 1$ 是等价无穷小，则 $a = \underline{\hspace{2cm}}$.
		
		\item 设 $\int_{1}^{x} f(t) \diff t = x [f(x) + 1]$，则 $f(x) = \underline{\hspace{2cm}}$.
		
		\item 对于函数 $f(x) = \ln \sin x$，在区间 $\left[ \frac{\pi}{4}, \frac{5\pi}{6} \right]$ 上使 Lagrange 中值定理结论成立的点是 $\xi = \underline{\hspace{2cm}}$.
		
		\item 已知函数 $f(x) = e^{-x} \ln ax$ 在 $x = \frac{1}{2}$ 处有极值，则 $a$ 的值为 \underline{\hspace{2cm}}.
		
		\item 设 $f(x) = \begin{cases} x + \pi, & -\pi \leq x < 0 \\ x - \pi, & 0 \leq x \leq \pi \end{cases}$，则 $f(x)$ 以 $2\pi$ 为周期的 Fourier 级数在 $[-\pi, \pi]$ 上的和函数为 $S(x) = \underline{\hspace{2cm}}$.
	\end{enumerate}
	
	\subsubsection*{二、选择题(本题15分，每小题3分)}
	\begin{enumerate}
		\item $\lim\limits_{x \to 0^{+}} \left( \ln \frac{1}{x} \right)^{x} = ( \quad )$
		\fourch{0}{1}{-1}{2}
		
		\item 当 $x > 0$ 时，曲线 $y = x \sin \frac{1}{x}$ ( \quad )
		\twoch{有且仅有水平渐近线}{有且仅有铅直渐近线}{既有水平渐近线也有铅直渐近线}{既无水平渐近线也无铅直渐近线}
		
		\item 下列反常积分中收敛的是 ( \quad )
		\fourch{$\int_{0}^{\infty} \cos x \diff x$}{$\int_{-1}^{1} \frac{\diff x}{x^2}$}{$\int_{0}^{1} \frac{\diff x}{\sqrt{x(1-x)}}$}{$\int_{-\infty}^{0} e^{-x} \diff x$}
		
		\item 双纽线 $(x^2 + y^2)^2 = x^2 - y^2$ 所围成的区域面积可用定积分表示为 ( \quad )
		\twoch{$2 \int_{0}^{\pi/4} \cos 2\theta d\theta$}{$2 \int_{0}^{\pi/2} \cos 2\theta d\theta$}{$4 \int_{0}^{\pi/4} \cos 2\theta d\theta$}{$\frac{1}{2} \int_{0}^{\pi/4} (\cos 2\theta)^2 d\theta$}
		
		\item 若级数 $\sum\limits_{n=1}^{\infty} (-1)^n \frac{(x-a)^n}{n}$ 在 $x > 0$ 时发散，在 $x = 0$ 时收敛，则常数 $a = ( \quad )$
		\fourch{-1}{1}{-2}{2}
	\end{enumerate}
	
	\subsubsection*{三、(12分)}
	设 $f(x)= \begin{cases} \frac{1-\cos ax}{x}, & x<0 \\ 0, & x=0, \\ \frac{\ln(b+x^2)}{x}, & x>0 \end{cases}$，试确定常数 $a,b$ 的值，使 $f(x)$ 在 $(-\infty,+\infty)$ 内处处可导.
	
	\subsubsection*{四、求解下列各题(共16分，每题8分)}
	\begin{enumerate}
		\item 求曲线 $e^{x+y}+xy=0$ 在点 $(1,-1)$ 处的切线与法线方程.
		
		\item 求曲线 $y=x^2-2x,y=0,x=1,x=3$ 所围成的平面图形绕y轴旋转一周所得的旋转体的体积V.
	\end{enumerate}
	
	\subsubsection*{五、(8分)}
	设某银行一年内吸纳储户存款的总数与银行付给储户年利率的平方成正比，若银行以20\%的年利率把储户存款总数的90\%贷出以获取利润，问银行支付给储户的年利率定为多少时，才能获得最大年利润？
	
	\subsubsection*{六、计算下列各题(共16分，每题8分)}
	\begin{enumerate}
		\item 计算不定积分： $\int \frac{\diff x}{e^x(e^x-1)}$.
		
		\item 设 $f(x)=\begin{cases} xe^{-x}, & x\leq0 \\ \sqrt{2x-x^2}, & 0<x\leq1 \end{cases}$，求 $\int_{-3}^1 f(x) \diff x$.
	\end{enumerate}
	
	\subsubsection*{七、(10分)}
	求幂级数 $\sum\limits_{n=1}^{\infty} \frac{x^n}{n\cdot 2^n}$ 的收敛域以及和函数 $S(x)$.
	
	\subsubsection*{八、(8分)}
	设 $f(x)$ 在 $[a,b]$ 上连续，在 $(a,b)$ 内可微，如果 $a\geq0$，证明：在 $(a,b)$ 内存在三个数 $x_1,x_2,x_3$，使 $f'(x_1)=(b+a)\frac{f'(x_2)}{2x_2}=(b^2+ba+a^2)\frac{f'(x_3)}{3x_3^2}$ 成立.
	
	\newpage
	\begin{spacing}{0}
		\begin{center}
			\textbf{\huge \kaishu 中南大学考试试卷}\\[10pt]
			\textbf{\large  \kaishu 2012\~{}2013学年\underline{一}学期\underline{高等数学A（一）}}\\[10pt]
			\kaishu(时间：2013年1月17日，星期四，10:00—11:40，共计：100分钟)\\[10pt]
			\underline{80}学时，\underline{5}学分，\underline{闭卷}，总分100分，占总评成绩\underline{60}\%
		\end{center}
	\end{spacing}
	
	\subsubsection*{一、填空题(本题15分，每小题3分)}
	\begin{enumerate}
		\item $d\int d\int d \int \frac{\sin x}{x} \diff x = \underline{\hspace{2cm}}$.
		
		\item 曲线 $y = \frac{x^2 + x}{x^2 - 1}$ 的渐近线的条数为 \underline{\hspace{2cm}}.
		
		\item $\int_0^{\infty} \frac{\ln x}{1 + x^2} \diff x = \underline{\hspace{2cm}}$.
		
		\item 设 $f(x) = \begin{cases} x + \pi, & -\pi \leq x < 0, \\ 0, & x = 0, \\ 1, & 0 < x \leq \pi \end{cases}$ 是以 $2\pi$ 为周期的 Fourier 级数的和函数为 $S(x)$，则 $S(x)$ 在 $[-\pi, \pi]$ 上的表达式为 \underline{\hspace{2cm}}.
		
		\item 设 $y = x(x-1)(x-2) \cdots (x-n)$，则 $y^{(n+1)} = \underline{\hspace{2cm}}$.
	\end{enumerate}
	
	\subsubsection*{二、选择题(本题15分，每小题3分)}
	\begin{enumerate}
		\item 当 $x \to 1$ 时，函数 $\frac{x^2 - 1}{x-1} e^{x-1}$ 的极限是（ ）
		\fourch{0}{2}{$\infty$}{不存在}
		
		\item 设函数 $f(x)$ 在 $(a,b)$ 内可导，$x_1, x_2 \in (a,b)$，则至少存在一点 $\xi$，使（ ）成立.
		\onech{$f(b) - f(a) = (b-a)f'(\xi), \xi \in (a,b)$}{$f(x_2) - f(x_1) = (x_2 - x_1)f'(\xi), \xi \in (a,b)$}{$f(x_2) - f(x_1) = (x_2 - x_1)f'(\xi), \xi$ 在 $x_1$ 与 $x_2$ 之间}{以上结论都不成立}
		
		\item 曲线 $y = x(x-1)(2-x)$ 与 $x$ 轴所围成的图形面积可表示为（ ）
		\onech{$-\int_0^2 x(x-1)(2-x)\diff x$}{$\int_0^1 x(x-1)(2-x)\diff x - \int_1^2 x(x-1)(2-x)\diff x$}{$\int_0^2 x(x-1)(2-x)\diff x$}{$-\int_0^1 x(x-1)(2-x)\diff x + \int_1^2 x(x-1)(2-x)\diff x$}
		
		\item 设 $\frac{d}{\diff x} f(x) = g(x), h(x) = x^2$，则 $\frac{d}{\diff x} f(h(x)) = $（）
		\fourch{$g(x^2)$}{$2xg(x)$}{$x^2 g(x^2)$}{$2xg(x^2)$}
		
		\item 若级数 $\sum\limits_{n=1}^{\infty} a_n^2$ 收敛，则级数 $\sum\limits_{n=1}^{\infty} a_n$ （ ）
		\fourch{一定绝对收敛}{一定条件收敛}{一定发散}{不能确定}
	\end{enumerate}
	
	\subsubsection*{三、（10分）}
	设 $f(x) = \begin{cases} a + x^2 & x < 0 \\ 1 & x = 0, \\ \ln(b + x^2) & x > 0 \end{cases}$，已知 $f(x)$ 在 $x = 0$ 处连续可导，求参数 $a, b$ 的值并求 $f'(x)$.
	
	\subsubsection*{四、（共15分，其中第一小题7分，第二小题8分）}
	\begin{enumerate}
		\item 求极限 $\lim\limits_{x \to 0} \frac{\arctan x^2 (1 - \cos \sqrt{x})}{\ln(1 + x)\sin^2 x}$.
		
		\item 设 $y = y(x)$ 由方程组 $\begin{cases} x = 3t^2 + 2t + 3 \\ e^y \sin t - y + 1 = 0 \end{cases}$ 所确立，求 $\frac{\diff y}{\diff x}|_{t=0}$.
	\end{enumerate}
	
	\subsubsection*{五、（10分）}
	求不定积分 $\int \frac{\sin x}{a \sin x + b \cos x} \diff x, \quad (a, b 不全为0)$.
	
	\subsubsection*{六、（共15分，其中第一小题7分，第二小题8分）}
	\begin{enumerate}
		\item 设 $f(x) = \int_0^x \frac{\sin t}{\pi - t} \diff t$，计算 $\int_0^\pi f(x) \diff x$.
		
		\item 问 $k$ 为何值时，由曲线 $y = x^2$，直线 $y = kx(0 < k < 2)$ 及 $x = 2$ 围成的平面图形面积最小？
	\end{enumerate}
	
	\subsubsection*{七、（12分）}
	设 $f(x) = \sum\limits_{n=1}^{\infty} n3^{n-1}x^{n-1}$，
	\begin{enumerate}
		\item 证明 $f(x)$ 在开区间 $\left(-\frac{1}{3}, \frac{1}{3}\right)$ 内连续；
		\item 求幂级数 $\sum\limits_{n=1}^{\infty} n3^{n-1}x^{n-1}$ 的收敛域；
		\item 计算定积分 $\int_0^{\frac{1}{6}} f(x) \diff x$ 的值.
	\end{enumerate}
	
	\subsubsection*{八、（8分）}
	设函数 $f(x)$ 在闭区间 $[0,4]$ 上连续，在开区间 $(0,4)$ 上可导，且 $f(0) + f(1) + f(2) + f(3) = 8, \quad f(4) = 2$，证明必存在 $x \in (0,4)$ 使 $f'(x) = 0$.
	
	
\end{document}
